\documentclass[11pt]{article}
\usepackage[a4paper,margin=20mm]{geometry}
\usepackage{amsmath,amssymb,bm,mathtools}
\usepackage{physics}
\usepackage{microtype}
\usepackage{booktabs}
\usepackage{enumitem}
\usepackage{xcolor}
\usepackage{hyperref}
\usepackage{titlesec}
\usepackage{slashed}
\usepackage[T1]{fontenc}
\usepackage{lmodern}
\hypersetup{colorlinks=true,linkcolor=blue,citecolor=blue,urlcolor=blue}
\titleformat{\section}{\large\bfseries}{\thesection}{0.7em}{}
\titleformat{\subsection}{\normalsize\bfseries}{\thesubsection}{0.7em}{}
\setlength{\parindent}{0pt}
\setlength{\parskip}{5pt}
\newcommand{\Id}{\mathbb I}
\newcommand{\Trc}{\operatorname{Tr}}
\newcommand{\ssb}{s\bar s}
\newcommand{\LLb}{\Lambda\bar\Lambda}
\newcommand{\eps}{\epsilon}
\newcommand{\Ccal}{\mathcal C}
\newcommand{\Dcal}{\mathcal D}
\newcommand{\Scal}{\mathcal S}
\newcommand{\Kcal}{\mathcal K}

\title{\textbf{Daily Research Progress Report -- 23 August 2026}\\
Microscopic transport of $s\bar s$ spin correlations toward $\Lambda\bar\Lambda$ observables}
\author{Detailed reconstruction of the full-day research program}
\date{23 August 2026}

\begin{document}
\maketitle

\begin{abstract}
This report reconstructs and expands the complete set of technical ideas developed on 23 August 2026 in the ongoing program to build a microscopic theory for the evolution of an initially produced strange-quark pair into the measured $\Lambda\bar\Lambda$ spin correlation. The central theme of the day was to replace a purely phenomenological ``decoherence factor'' by a hierarchy of increasingly microscopic descriptions. We begin with the exact two-spin density matrix and its connected correlation tensor, derive the two-body BBGKY equation, and show explicitly where closure approximations enter. We then derive the relativistically correct projection of the axial Wigner tensor onto rest-frame spin components, demonstrating why a naive spatial trace in the laboratory frame is not generally the measured scalar. For Markovian one-body relaxation we prove that the two-spin mode decays with the sum of the two leg rates. For stochastic background fields we derive, rather than assume, the cancellation of common-mode fluctuations and obtain the relative-noise rate from a double-commutator master equation. We extend this result to finite-memory noise with a second-order time-convolutionless kernel and show the crossover from quadratic short-time decay to exponential long-time decay. We then construct the gauge-covariant QCD interpretation in terms of Wilson-line-connected field-strength correlators, clarify which parts of this language are controlled and which remain model dependent for strange quarks, and connect the result to the Kadanoff--Baym/Wigner program. The latter is developed in detail through the Wigner transform, Moyal product, Clifford decomposition, and the need for a connected four-point function for pair spin correlations. Finally we discuss source-propagator separation, hadronization and feed-down, positivity constraints, and the qualitatively different STAR and CMS phenomenology. The report ends with a prioritized derivation program for obtaining a genuine QCD singlet-triplet and tensor collision kernel.
\end{abstract}

\tableofcontents
\newpage

\section{Research questions attacked during the day}

The six research sessions on 23 August were not independent topics. They formed a sequence of questions intended to move the project from a hadron-level fit toward a microscopic nonequilibrium-QCD description. The questions can be organized as follows.

\begin{enumerate}[leftmargin=1.7em]
\item What is the minimal density-matrix object that contains the spin information measured by STAR and CMS, and when must one use the \emph{connected} rather than the full two-spin tensor?
\item Starting from an exact many-body density operator, what is the exact equation of motion for the reduced $s\bar s$ density matrix, and where does a kinetic closure enter?
\item How should relativistic quark spin be projected from the axial Wigner function to the three-component spin tensor that is eventually mapped to hyperon decay angles?
\item If each strange constituent undergoes independent spin relaxation, how does the pair correlation decay? In particular, do relaxation rates add, multiply, or interfere?
\item For a common fluctuating color environment, which part of the environmental field destroys the scalar correlation? Can perfectly common noise decohere the STAR scalar?
\item How is the white-noise result modified when the environment has a finite memory time?
\item What is the gauge-covariant QCD object underlying the effective cross-noise kernel, and why is a Wilson line required?
\item How can the phenomenological master-equation language be embedded in a Kadanoff--Baym/Wigner derivation without double counting the source and the subsequent medium evolution?
\item What pieces of hadronization, feed-down, and source mixing remain indispensable before a partonic calculation can be compared with the observed $\Lambda\bar\Lambda$ quantity?
\item What can already be learned from the qualitative difference between the STAR result at $\sqrt{s}=200$ GeV and the CMS result at LHC energies?
\end{enumerate}

The working philosophy throughout this report is to distinguish four levels of statement:

\begin{itemize}[leftmargin=1.5em]
\item \textbf{Exact identities}, following only from quantum mechanics and definitions;
\item \textbf{Controlled approximations}, such as a gradient expansion or a weak-coupling/Markov limit;
\item \textbf{Model assumptions}, for example Gaussian stochastic color fields;
\item \textbf{Phenomenological factorization ansatze}, useful for fitting data but not yet derived from QCD.
\end{itemize}

This distinction is essential because the long-term goal is not merely to reproduce one curve but to identify which microscopic QCD structures the data actually constrain.

\section{Experimental target and density-matrix conventions}

\subsection{General two-spin density matrix}

For two spin-$1/2$ degrees of freedom, any $4\times4$ Hermitian density matrix can be expanded in the Pauli basis. Since
\begin{equation}
\{\Id,\sigma_x,\sigma_y,\sigma_z\}
\end{equation}
forms an orthogonal operator basis for a single qubit, the tensor products form a basis for the two-spin system. We therefore write
\begin{equation}
\rho_{12}
=
\frac14\left[
\Id\otimes\Id
+p_i\sigma_i\otimes\Id
+\bar p_j\Id\otimes\sigma_j
+C^{\rm full}_{ij}\sigma_i\otimes\sigma_j
\right].
\label{eq:rho_general}
\end{equation}
Repeated Cartesian indices are summed.

The factor $1/4$ is fixed by normalization. Indeed,
\begin{align}
\Trc\rho_{12}
&=\frac14\Trc(\Id\otimes\Id)\\
&=\frac14(2)(2)=1,
\end{align}
because every term containing one Pauli matrix has zero trace.

Using
\begin{equation}
\Trc(\sigma_i\sigma_j)=2\delta_{ij},
\end{equation}
we recover the coefficients directly:
\begin{align}
p_i&=\Trc[\rho_{12}(\sigma_i\otimes\Id)],\\
\bar p_j&=\Trc[\rho_{12}(\Id\otimes\sigma_j)],\\
C^{\rm full}_{ij}&=\Trc[\rho_{12}(\sigma_i\otimes\sigma_j)].
\end{align}

\subsection{Full versus connected correlation tensor}

The quantity $C^{\rm full}_{ij}$ contains both genuine pair correlation and the trivial product of one-body polarizations. The connected tensor is
\begin{equation}
\boxed{
C_{ij}
\equiv
C^{\rm full}_{ij}-p_i\bar p_j
=
\langle\sigma_{1i}\sigma_{2j}\rangle
-
\langle\sigma_{1i}\rangle
\langle\sigma_{2j}\rangle.
}
\label{eq:Cconnected}
\end{equation}

Why is this distinction important? Suppose the two spins are completely uncorrelated,
\begin{equation}
\rho_{12}=\rho_1\otimes\rho_2,
\end{equation}
with
\begin{equation}
\rho_1=\frac12(\Id+p_i\sigma_i),
\qquad
\rho_2=\frac12(\Id+\bar p_j\sigma_j).
\end{equation}
Then
\begin{equation}
C^{\rm full}_{ij}=p_i\bar p_j,
\end{equation}
while
\begin{equation}
C_{ij}=0.
\end{equation}
Thus a theory of pair correlation should propagate the connected object whenever nonzero one-body polarization is possible.

In the STAR $pp$ application, the pair-averaged one-body polarizations are small enough that the full and connected scalar are often numerically close. Nevertheless, the distinction must remain explicit in a general heavy-ion extension where global or local hyperon polarization can be nonzero.

\subsection{Scalar projection}

Define
\begin{equation}
Q\equiv\bm\sigma_1\cdot\bm\sigma_2=\sum_{a=1}^{3}\sigma_{1a}\sigma_{2a}.
\end{equation}
For vanishing one-body polarizations, the scalar relative-polarization observable can be written
\begin{equation}
\boxed{
P=\frac13\langle Q\rangle
=\frac13\delta_{ab}C_{ab}.
}
\label{eq:Pscalar}
\end{equation}

For nonzero one-body polarization the connected scalar is instead
\begin{equation}
P_{\rm conn}
=\frac13\delta_{ab}C_{ab}
=\frac13\left[
\langle Q\rangle-\bm p\cdot\bar{\bm p}
\right].
\end{equation}

The STAR short-range result is
\begin{equation}
P_{\Lambda\bar\Lambda}^{\rm STAR}
=0.181\pm0.035_{\rm stat}\pm0.022_{\rm sys},
\label{eq:STARnum}
\end{equation}
with the signal decreasing strongly at large pair separation. At the LHC, CMS finds no statistically significant $\Lambda\bar\Lambda$ deviation from zero within uncertainties, while same-sign hyperon pairs show negative correlations at small $\Delta R$. These two experimental facts already warn us that the source and/or propagation mechanism is not universal across collision energies.

\section{Exact two-body reduction from the many-body density operator}

\subsection{Many-body equation}

Let the full $N$-body density operator obey
\begin{equation}
i\partial_t\rho_N=[H_N,\rho_N].
\label{eq:vonN}
\end{equation}
Take a Hamiltonian with one-body terms $h_i$ and pair interactions $V_{ij}$,
\begin{equation}
H_N=\sum_{i=1}^{N}h_i+\frac12\sum_{i\neq j}V_{ij}.
\end{equation}
The two-body reduced density operator is
\begin{equation}
\rho_{12}
=\Trc_{3\cdots N}\rho_N.
\end{equation}

Differentiate and use Eq.~\eqref{eq:vonN}:
\begin{align}
i\partial_t\rho_{12}
&=\Trc_{3\cdots N}[H_N,\rho_N].
\end{align}
Terms involving only particles $1$ and $2$ commute through the spectator trace:
\begin{equation}
\Trc_{3\cdots N}[h_1+h_2+V_{12},\rho_N]
=[h_1+h_2+V_{12},\rho_{12}].
\end{equation}

Terms containing only spectators disappear after the partial trace because of cyclicity within the spectator Hilbert space. The nontrivial terms couple particle $1$ or $2$ to a spectator. For particle $3$, for example,
\begin{equation}
\Trc_3[V_{13}+V_{23},\rho_{123}].
\end{equation}
Summing the spectators gives the second BBGKY equation,
\begin{equation}
\boxed{
i\partial_t\rho_{12}
=[h_1+h_2+V_{12},\rho_{12}]
+\sum_{k=3}^{N}\Trc_k[V_{1k}+V_{2k},\rho_{12k}].
}
\label{eq:BBGKY2}
\end{equation}

\subsection{Where the closure problem enters}

Equation~\eqref{eq:BBGKY2} is exact but not closed because it contains the three-body reduced density operator. Write a cluster expansion schematically as
\begin{equation}
\rho_{123}
=\rho_1\rho_2\rho_3
+g_{12}\rho_3+g_{13}\rho_2+g_{23}\rho_1
+g_{123},
\end{equation}
where $g_{ij}$ are connected pair correlations and $g_{123}$ is the genuine three-body cumulant.

A molecular-chaos or Born approximation discards $g_{123}$ and approximates the spectator as a bath state. A kinetic equation then emerges only after assumptions about

\begin{enumerate}[leftmargin=1.6em]
\item the time scale on which bath correlations decay,
\item whether the bath is weakly perturbed by the pair,
\item whether memory effects can be neglected,
\item and how color gauge covariance is maintained.
\end{enumerate}

This derivation is conceptually useful because it identifies exactly what a phenomenological master equation is replacing: the spectator trace in Eq.~\eqref{eq:BBGKY2}, together with the hierarchy of higher correlations.

\subsection{Equation for the pair tensor}

For any operator $O_{12}$,
\begin{equation}
\frac{d}{dt}\langle O_{12}\rangle
=-i\langle[O_{12},H_{12}]\rangle
-i\sum_{k\ge3}\Trc_{12k}\left[O_{12}[V_{1k}+V_{2k},\rho_{12k}]\right].
\end{equation}
Choosing
\begin{equation}
O_{12}=\sigma_i\otimes\sigma_j
\end{equation}
shows that the pair spin tensor is coupled to mixed spin-color-space correlations with spectators. This is the microscopic origin of both coherent mean fields and dissipative collision kernels.

\section{Relativistic spin projection from Dirac bilinears}

\subsection{Spin four-vector}

For a positive-energy Dirac spinor with four-momentum $p^\mu$ and spin four-vector $S^\mu$, the covariant spin projector is
\begin{equation}
u(p,S)\bar u(p,S)
=\frac12(\slashed p+m)(1+\gamma^5\slashed S),
\label{eq:spinprojector}
\end{equation}
with
\begin{equation}
p\cdot S=0,
\qquad
S^2=-1.
\end{equation}

To verify the axial-vector bilinear, use Eq.~\eqref{eq:spinprojector}:
\begin{align}
\bar u\gamma^\mu\gamma^5u
&=\Trc\left[\gamma^\mu\gamma^5u\bar u\right]\\
&=\frac12\Trc\left[\gamma^\mu\gamma^5(\slashed p+m)(1+\gamma^5\slashed S)\right].
\end{align}
The terms without $S$ vanish because they contain traces with an odd number of gamma matrices or a single $\gamma^5$. The surviving piece is
\begin{align}
\frac12\Trc\left[\gamma^\mu\gamma^5(\slashed p+m)\gamma^5\slashed S\right].
\end{align}
Using $\gamma^5\slashed p\gamma^5=-\slashed p$ and $\gamma^5m\gamma^5=m$, the term with $\slashed p\slashed S$ contributes a trace of three gamma matrices and vanishes, while
\begin{align}
\frac12 m\Trc(\gamma^\mu\slashed S)
&=\frac12m(4S^\mu)\\
&=2mS^\mu.
\end{align}
Thus
\begin{equation}
\boxed{
\bar u(p,S)\gamma^\mu\gamma^5u(p,S)=2mS^\mu.
}
\label{eq:axialbilinear}
\end{equation}

This identity explains why the axial component of the Wigner function is the natural relativistic carrier of spin information.

\subsection{Momentum-dependent spin tetrad}

A three-component spin vector is not a Lorentz four-vector with three arbitrary spatial components. For each momentum $p^\mu$ choose three spacelike basis vectors $e_a^\mu(p)$ satisfying
\begin{equation}
e_a\cdot p=0,
\qquad
e_a\cdot e_b=-\delta_{ab}.
\label{eq:tetrad}
\end{equation}
Together with $p^\mu/m$, these form an orthonormal tetrad.

The spin four-vector is expanded as
\begin{equation}
S^\mu=s_a e_a^\mu(p).
\end{equation}
Because of Eq.~\eqref{eq:tetrad},
\begin{equation}
s_a=-e_{a\mu}S^\mu.
\end{equation}
The overall sign depends on the metric and basis convention; what matters is that the same convention is used on both legs.

\subsection{Two-particle axial tensor}

Introduce a connected axial--axial tensor in phase space,
\begin{equation}
\Ccal^{\mu\nu}(1,2)
\equiv
\langle\mathcal A_s^\mu(1)\mathcal A_{\bar s}^\nu(2)\rangle_c.
\end{equation}
The three-dimensional spin tensor is obtained by projecting each Lorentz index onto its own momentum-dependent spin basis:
\begin{equation}
\boxed{
C_{ab}(1,2)
=e_{1a\mu}(p_1)e_{2b\nu}(p_2)\Ccal^{\mu\nu}(1,2),
}
\label{eq:tetradprojection}
\end{equation}
up to normalization factors associated with the convention used for $\mathcal A^\mu$.

Consequently the scalar is
\begin{equation}
\boxed{
P(1,2)=\frac13\sum_{a=1}^{3}e_{1a\mu}e_{2a\nu}\Ccal^{\mu\nu}(1,2).
}
\label{eq:relP}
\end{equation}

A naive laboratory-frame quantity
\begin{equation}
\frac13\sum_{i=x,y,z}\Ccal^{ii}
\end{equation}
is not generally equivalent to Eq.~\eqref{eq:relP}, because the physical spin bases depend on $p_1$ and $p_2$. This is a genuine relativistic correction, not a matter of notation.

\section{Independent one-body spin relaxation}

\subsection{Single-spin depolarizing generator}

Consider one spin with an isotropic depolarizing generator
\begin{equation}
\Dcal_m\rho
=\frac{\gamma_m}{4}\sum_{a=1}^{3}
(\sigma_{ma}\rho\sigma_{ma}-\rho).
\label{eq:depgen}
\end{equation}
We derive its action on a Pauli operator in the Heisenberg picture. The adjoint is
\begin{equation}
\Dcal_m^\dagger O
=\frac{\gamma_m}{4}\sum_a
(\sigma_{ma}O\sigma_{ma}-O).
\end{equation}
For $O=\sigma_{mi}$, use
\begin{equation}
\sum_a\sigma_a\sigma_i\sigma_a=-\sigma_i.
\label{eq:pauliidentity}
\end{equation}
To check Eq.~\eqref{eq:pauliidentity}, the term $a=i$ gives $\sigma_i$, while each of the other two Pauli matrices anticommutes with $\sigma_i$ and contributes $-\sigma_i$. Hence
\begin{equation}
\sigma_i-\sigma_i-\sigma_i=-\sigma_i.
\end{equation}
Therefore
\begin{align}
\Dcal_m^\dagger\sigma_{mi}
&=\frac{\gamma_m}{4}\left[-\sigma_{mi}-3\sigma_{mi}\right]\\
&=-\gamma_m\sigma_{mi}.
\end{align}
So the one-body polarization satisfies
\begin{equation}
\dot p_i=-\gamma_m p_i.
\end{equation}

\subsection{Two-spin tensor under independent channels}

Let
\begin{equation}
\dot\rho=(\Dcal_1+\Dcal_2)\rho.
\end{equation}
For the product operator $\sigma_i\otimes\sigma_j$,
\begin{align}
(\Dcal_1^\dagger+\Dcal_2^\dagger)(\sigma_i\otimes\sigma_j)
&=(-\gamma_1\sigma_i)\otimes\sigma_j
+\sigma_i\otimes(-\gamma_2\sigma_j)\\
&=-(\gamma_1+\gamma_2)\sigma_i\otimes\sigma_j.
\end{align}
Hence
\begin{equation}
\boxed{
\dot C_{ij}=-(\gamma_s+\gamma_{\bar s})C_{ij}
}
\label{eq:indepC}
\end{equation}
when connected one-body terms are treated consistently.

For time-dependent rates,
\begin{equation}
\boxed{
C_{ij}(t_h)
=C_{ij}(t_0)
\exp\left[-\int_{t_0}^{t_h}dt\,(\gamma_s+\gamma_{\bar s})\right].
}
\label{eq:indepSolution}
\end{equation}
The scalar obeys the same multiplicative factor.

This provides a useful consistency rule: if the two legs lose spin information independently, their exponents add. A phenomenological formula using only one leg's rate would undercount the pair damping by a factor of two in the symmetric case.

\subsection{Interpretation of perturbative strange-quark spin equilibration}

Perturbative estimates of strange-quark spin equilibration in a weakly coupled rotating QGP find times much longer than typical heavy-ion evolution times. This does not prove that spin correlations are frozen during confinement, because the STAR phenomenon occurs in a nonperturbative hadronization regime rather than a weakly coupled thermal QGP. It does, however, imply that ordinary perturbative helicity/spin-flip scattering is unlikely to be the dominant mechanism responsible for a rapid loss of the pair correlation.

\section{Common and relative stochastic fields}

\subsection{Hamiltonian decomposition}

Model unresolved spin-dependent environmental fields by
\begin{equation}
H_I(t)
=\frac12\left[
\bm\xi_1(t)\cdot\bm\sigma_1
+\bm\xi_2(t)\cdot\bm\sigma_2
\right].
\label{eq:stochH}
\end{equation}
Define the common and relative fields
\begin{equation}
\bar{\bm\xi}
=\frac{\bm\xi_1+\bm\xi_2}{2},
\qquad
\bm\xi_-
=\bm\xi_1-\bm\xi_2.
\end{equation}
Then
\begin{equation}
H_I
=\frac12\bar{\bm\xi}\cdot(\bm\sigma_1+\bm\sigma_2)
+\frac14\bm\xi_-\cdot(\bm\sigma_1-\bm\sigma_2).
\label{eq:commonrelative}
\end{equation}

\subsection{Exact cancellation of common-mode precession}

Recall
\begin{equation}
Q=\sum_i\sigma_{1i}\sigma_{2i}.
\end{equation}
We compute
\begin{align}
[Q,\sigma_{1a}]
&=\sum_i[\sigma_{1i},\sigma_{1a}]\sigma_{2i}\\
&=2i\sum_{ik}\eps_{iak}\sigma_{1k}\sigma_{2i},
\end{align}
and similarly
\begin{align}
[Q,\sigma_{2a}]
&=2i\sum_{ik}\eps_{iak}\sigma_{1i}\sigma_{2k}.
\end{align}
Rename $i\leftrightarrow k$ in the second term:
\begin{equation}
[Q,\sigma_{2a}]
=2i\sum_{ik}\eps_{kai}\sigma_{1k}\sigma_{2i}.
\end{equation}
Because $\eps_{kai}=-\eps_{iak}$,
\begin{equation}
[Q,\sigma_{2a}]=-[Q,\sigma_{1a}].
\end{equation}
Therefore
\begin{equation}
\boxed{
[Q,\sigma_{1a}+\sigma_{2a}]=0.
}
\label{eq:commoncancel}
\end{equation}

Equation~\eqref{eq:commoncancel} is exact. A perfectly common field generates a collective $SU(2)$ rotation of the pair and leaves $Q$ invariant. Thus the scalar correlation is not sensitive to the magnitude of common spin precession; it is sensitive to the difference between what the two constituents experience.

This is one of the most robust results in the whole research program.

\section{White-noise limit and the relative-noise damping rate}

\subsection{Gaussian noise correlator}

Assume zero-mean Gaussian noise
\begin{equation}
\langle\xi_m^a(t)\rangle=0,
\end{equation}
with
\begin{equation}
\langle\xi_m^a(t)\xi_n^b(t')\rangle
=2D_{mn}\delta^{ab}\delta(t-t').
\label{eq:whitenoise}
\end{equation}
Here $m,n\in\{1,2\}$ label the two particles. Positivity of the covariance matrix requires
\begin{equation}
D_{11}\ge0,
\qquad
D_{22}\ge0,
\qquad
D_{12}^2\le D_{11}D_{22}.
\label{eq:Dpositive}
\end{equation}

For Hermitian classical Gaussian noise, averaging the stochastic unitary evolution gives the double-commutator master equation
\begin{equation}
\boxed{
\dot\rho
=-\frac14\sum_{mn,a}D_{mn}
[\sigma_m^a,[\sigma_n^a,\rho]].
}
\label{eq:doublecomm}
\end{equation}
The overall normalization follows from the $1/2$ in Eq.~\eqref{eq:stochH} and the convention in Eq.~\eqref{eq:whitenoise}.

\subsection{Equation for the full tensor}
\begin{equation}
O_{ij}=\sigma_{1i}\sigma_{2j}.
\end{equation}
Using the adjoint form,
\begin{equation}
\frac{d}{dt}\langle O_{ij}\rangle
=-\frac14\sum_{mn,a}D_{mn}
\langle[\sigma_n^a,[\sigma_m^a,O_{ij}]]\rangle.
\end{equation}

For $m=n=1$,
\begin{align}
[\sigma_1^a,O_{ij}]
&=[\sigma_1^a,\sigma_{1i}]\sigma_{2j}\\
&=2i\eps_{aik}\sigma_{1k}\sigma_{2j}.
\end{align}
A second commutator gives
\begin{align}
[\sigma_1^a,[\sigma_1^a,O_{ij}]]
&=2i\eps_{aik}[\sigma_1^a,\sigma_{1k}]\sigma_{2j}\\
&=-4\eps_{aik}\eps_{ak\ell}\sigma_{1\ell}\sigma_{2j}.
\end{align}
Summing over $a,k$ gives
\begin{equation}
\sum_{ak}\eps_{aik}\eps_{ak\ell}=-2\delta_{i\ell},
\end{equation}
so
\begin{equation}
\sum_a[\sigma_1^a,[\sigma_1^a,O_{ij}]]=8O_{ij}.
\end{equation}
Thus the $D_{11}$ piece contributes $-2D_{11}C_{ij}$, and similarly the $D_{22}$ piece gives $-2D_{22}C_{ij}$.

For the cross term,
\begin{align}
[\sigma_1^a,O_{ij}]&=2i\eps_{aik}\sigma_{1k}\sigma_{2j},\\
[\sigma_2^a,[\sigma_1^a,O_{ij}]]
&=(2i)^2\eps_{aik}\eps_{aj\ell}\sigma_{1k}\sigma_{2\ell}\\
&=-4\eps_{aik}\eps_{aj\ell}\sigma_{1k}\sigma_{2\ell}.
\end{align}
Use
\begin{equation}
\sum_a\eps_{aik}\eps_{aj\ell}
=\delta_{ij}\delta_{k\ell}-\delta_{i\ell}\delta_{kj}.
\end{equation}
Hence
\begin{equation}
\sum_a[\sigma_2^a,[\sigma_1^a,O_{ij}]]
=-4(\delta_{ij}O_{kk}-O_{ji}).
\end{equation}
Accounting for both $D_{12}$ and $D_{21}$ and taking a symmetric covariance matrix gives
\begin{equation}
\boxed{
\dot C_{ij}
=-2(D_{11}+D_{22})C_{ij}
+2D_{12}(\delta_{ij}C_{kk}-C_{ji}).
}
\label{eq:Cwhite}
\end{equation}

\subsection{Isotropic scalar limit}

For an isotropic tensor
\begin{equation}
C_{ij}=P\delta_{ij},
\end{equation}
we have
\begin{equation}
C_{kk}=3P,
\qquad
C_{ji}=P\delta_{ij}.
\end{equation}
Substitute into Eq.~\eqref{eq:Cwhite}:
\begin{align}
\dot P\delta_{ij}
&=-2(D_{11}+D_{22})P\delta_{ij}
+2D_{12}(3P\delta_{ij}-P\delta_{ij})\\
&=\left[-2(D_{11}+D_{22})+4D_{12}\right]P\delta_{ij}.
\end{align}
Therefore
\begin{equation}
\boxed{
\dot P
=-2(D_{11}+D_{22}-2D_{12})P.
}
\label{eq:Pwhitegeneral}
\end{equation}

For statistically equivalent trajectories,
\begin{equation}
D_{11}=D_{22}=D(0),
\qquad
D_{12}=D(r),
\end{equation}
so
\begin{equation}
\boxed{
\Gamma_{\rm rel}(r)=4[D(0)-D(r)],
\qquad
\dot P=-\Gamma_{\rm rel}(r)P.
}
\label{eq:Gammarel}
\end{equation}

Several checks are immediate:

\begin{itemize}[leftmargin=1.5em]
\item At $r=0$, perfectly common noise has $D(r)=D(0)$ and $\Gamma_{\rm rel}=0$.
\item At large separation, $D(r)\to0$, giving independent-leg damping $\Gamma_{\rm rel}\to4D(0)$.
\item Positivity of the covariance ensures $D(r)\le D(0)$ for equivalent variables, so the rate is nonnegative.
\end{itemize}

\section{Finite-memory noise and non-Markovian decay}

\subsection{Interaction-picture expansion}

White noise is a strong assumption. Let the interaction-picture Hamiltonian be
\begin{equation}
H_I(t)=\frac12\sum_{m,a}\xi_m^a(t)\sigma_m^a.
\end{equation}
For zero-mean noise the first cumulant vanishes. At second order, the time-convolutionless equation has the form
\begin{equation}
\dot\rho(t)
=-\int_0^t d\tau\,
\langle[H_I(t),[H_I(t-\tau),\rho(t)]]\rangle
+\mathcal O(\xi^3),
\label{eq:TCL2}
\end{equation}
where the replacement $\rho(t-\tau)\to\rho(t)$ is the defining time-local TCL2 step. It does not yet require a delta-correlated bath.

Assume
\begin{equation}
\langle\xi_m^a(t)\xi_n^b(t-\tau)\rangle
=\delta^{ab}K_{mn}(\tau).
\label{eq:memorycorr}
\end{equation}
Repeating the same Pauli algebra as in the white-noise case, the isotropic scalar obeys
\begin{equation}
\boxed{
\dot P(t)
=-2P(t)\int_0^t d\tau
\left[K_{11}(\tau)+K_{22}(\tau)-K_{12}(\tau)-K_{21}(\tau)\right].
}
\label{eq:TCLP}
\end{equation}
For equivalent trajectories and a symmetric correlator,
\begin{equation}
\boxed{
\dot P(t)
=-4P(t)\int_0^t d\tau\,[K_0(\tau)-K_r(\tau)].
}
\label{eq:TCLPsym}
\end{equation}

The same relative-noise structure survives: memory modifies the time dependence, but common-mode cancellation remains.

\subsection{Exponential memory kernel}

Take
\begin{equation}
K_0(\tau)-K_r(\tau)
=K_\Delta e^{-\tau/\tau_c}.
\label{eq:expkernel}
\end{equation}
Then
\begin{align}
\int_0^t d\tau\,[K_0-K_r]
&=K_\Delta\tau_c(1-e^{-t/\tau_c}).
\end{align}
Equation~\eqref{eq:TCLPsym} becomes
\begin{equation}
\frac{\dot P}{P}
=-4K_\Delta\tau_c(1-e^{-t/\tau_c}).
\end{equation}
Integrate from $0$ to $t$:
\begin{align}
\ln\frac{P(t)}{P(0)}
&=-4K_\Delta\tau_c
\int_0^t dt'\,(1-e^{-t'/\tau_c})\\
&=-4K_\Delta\tau_c
\left[t-\tau_c(1-e^{-t/\tau_c})\right].
\end{align}
Thus
\begin{equation}
\boxed{
P(t)=P(0)
\exp\left\{
-4K_\Delta\tau_c
\left[t-\tau_c(1-e^{-t/\tau_c})\right]
\right\}.
}
\label{eq:memorysolution}
\end{equation}

\subsection{Short- and long-time limits}

For $t\ll\tau_c$,
\begin{equation}
1-e^{-t/\tau_c}
=\frac{t}{\tau_c}-\frac{t^2}{2\tau_c^2}+\cdots,
\end{equation}
so
\begin{align}
t-\tau_c(1-e^{-t/\tau_c})
&=\frac{t^2}{2\tau_c}+\mathcal O(t^3).
\end{align}
Therefore
\begin{equation}
P(t)
\simeq
P(0)e^{-2K_\Delta t^2}.
\label{eq:quadratic}
\end{equation}
The initial decay is quadratic, as expected for a unitary system observed before a Markovian coarse graining is valid.

For $t\gg\tau_c$,
\begin{equation}
t-\tau_c(1-e^{-t/\tau_c})\to t-\tau_c,
\end{equation}
so the asymptotic rate is
\begin{equation}
\Gamma_{\infty}=4K_\Delta\tau_c.
\end{equation}
This reproduces the exponential Markovian behavior after the bath memory has been integrated out.

\section{Spatial structure and a QCD correlation length}

Let the equal-time cross correlator decrease with pair separation,
\begin{equation}
D(r)=D(0)f(r/\ell_c),
\qquad
f(0)=1,
\qquad
f(\infty)=0.
\end{equation}
Then
\begin{equation}
\Gamma_{\rm rel}(r)=4D(0)[1-f(r/\ell_c)].
\end{equation}
For a smooth Gaussian choice
\begin{equation}
f(x)=e^{-x^2/2},
\end{equation}
we have at small separation
\begin{equation}
1-f(r/\ell_c)
=\frac{r^2}{2\ell_c^2}+\mathcal O(r^4),
\end{equation}
so the decoherence rate starts quadratically in $r$. For a nonanalytic exponential correlator
\begin{equation}
f(x)=e^{-|x|},
\end{equation}
the leading behavior is linear in $r$.

This gives a physically meaningful target for high-statistics data: the curvature of the small-separation spin-correlation curve is sensitive to the short-distance structure of the environmental correlator, not only to one overall damping constant.

However, the experimental variable is $\Delta R$ in momentum space, not a direct coordinate-space separation. A model is required to relate
\begin{equation}
r(t)\leftrightarrow \Delta R,
\end{equation}
for example through parton trajectories during hadronization. This is one reason an event generator or transport model is ultimately needed.

\section{Gauge-covariant QCD interpretation}

\subsection{Why a naive color-field correlator is not gauge invariant}

A schematic chromomagnetic correlator
\begin{equation}
\langle B^{ai}(x_1)B^{aj}(x_2)\rangle
\end{equation}
is not gauge invariant because the color field at $x_1$ and at $x_2$ transform in different local color frames. To compare them, parallel transport is required.

In the adjoint representation, introduce a Wilson line
\begin{equation}
\mathcal U^{ab}_{C}(x_1,x_2)
=\left[\mathcal P\exp\left(ig\int_C dz^\mu A_\mu^c(z)T^c_{\rm adj}\right)\right]^{ab}.
\end{equation}
Then a gauge-covariant two-point object is
\begin{equation}
\boxed{
\mathcal G^{ij}(x_1,x_2)
=
\left\langle
B^{ai}(x_1)
\mathcal U^{ab}_{C}(x_1,x_2)
B^{bj}(x_2)
\right\rangle.
}
\label{eq:gaugecorr}
\end{equation}
The path $C$ is part of the definition. In a dynamical pair problem the natural path may be related to the quark trajectories or to a Schwinger-Keldysh contour construction.

\subsection{Effective spin coupling and the strange-quark caveat}

For a heavy quark, a Pauli chromomagnetic interaction appears systematically in a $1/m_Q$ expansion. For the strange quark, $m_s$ is not parametrically larger than $\Lambda_{\rm QCD}$, so importing the heavy-quark coefficient as a controlled approximation is not justified.

The correct strategy is therefore:

\begin{enumerate}[leftmargin=1.6em]
\item use the field correlator structure, Eq.~\eqref{eq:gaugecorr}, as the symmetry-guided object;
\item derive the coefficient and Lorentz structure from relativistic axial kinetic theory or a microscopic model;
\item only then identify the effective $D_{mn}$ or $K_{mn}(\tau)$ that appears in the reduced spin equation.
\end{enumerate}

This separation prevents a conceptual mistake: a mathematically consistent Lindblad model can still have quantitatively incorrect QCD coefficients if its microscopic matching is uncontrolled.

\section{Unitality and what relative noise can and cannot do}

The Gaussian Hermitian-noise channel in Eq.~\eqref{eq:doublecomm} is unital:
\begin{equation}
\mathcal L(\Id)=0.
\end{equation}
Indeed every double commutator with the identity vanishes. Consequently
\begin{equation}
\rho=\frac{\Id_4}{4}
\end{equation}
is a fixed point.

This means relative noise can \emph{destroy or reshape an existing spin correlation}, but cannot generate a pair correlation starting from the completely mixed state. A genuine source of correlation must therefore belong to a different microscopic sector, for example

\begin{itemize}[leftmargin=1.5em]
\item correlated pair production,
\item a coherent pair interaction acting on a nontrivial source state,
\item a non-unital dissipative process,
\item or a connected collision/source term in the four-point kinetic hierarchy.
\end{itemize}

This distinction is important when combining glasma calculations with hadronization decoherence: a field-induced pair correlation generated at production must not be double counted as an additive source inside a later unital depolarizing channel.

\section{Kadanoff--Baym equations and Wigner transformation}

\subsection{One-particle lesser Green function}

Define the lesser fermion Green function
\begin{equation}
G^<(x_1,x_2)
=i\langle\bar\psi(x_2)\psi(x_1)\rangle
\end{equation}
(up to the sign convention used for fermions). Dyson equations on the Schwinger-Keldysh contour lead to left- and right-acting Kadanoff--Baym equations. Schematically,
\begin{equation}
(i\slashed\partial_{x_1}-m)G^<
=\Sigma^R\circ G^<+\Sigma^<\circ G^A,
\label{eq:KBleft}
\end{equation}
and
\begin{equation}
G^<(-i\overleftarrow{\slashed\partial}_{x_2}-m)
=G^R\circ\Sigma^<+G^<\circ\Sigma^A.
\label{eq:KBright}
\end{equation}
The symbol $\circ$ denotes a space-time convolution.

The sum and difference of Eqs.~\eqref{eq:KBleft} and \eqref{eq:KBright} generate constraint and transport equations after Wigner transformation.

\subsection{Wigner coordinates}

Introduce
\begin{equation}
X=\frac{x_1+x_2}{2},
\qquad
s=x_1-x_2,
\end{equation}
and define
\begin{equation}
G^<(X,p)
=\int d^4s\,e^{ip\cdot s}
G^<\left(X+\frac{s}{2},X-\frac{s}{2}\right).
\label{eq:wigner}
\end{equation}
Derivatives transform according to
\begin{equation}
\partial_{x_1}
=\frac12\partial_X+\partial_s,
\qquad
\partial_{x_2}
=\frac12\partial_X-\partial_s.
\end{equation}
Under the Fourier transform,
\begin{equation}
\partial_s\to -ip,
\end{equation}
so drift and mass-shell structures naturally emerge.

\subsection{Moyal product}

A convolution becomes a star product. For two Wigner-transformed quantities,
\begin{equation}
(A\circ B)(X,p)\to A\star B,
\end{equation}
with
\begin{equation}
\boxed{
A\star B
=A\exp\left[
\frac{i}{2}
\left(
\overleftarrow\partial_X\cdot\overrightarrow\partial_p
-
\overleftarrow\partial_p\cdot\overrightarrow\partial_X
\right)
\right]B.
}
\label{eq:moyal}
\end{equation}
Expanding to first order in gradients gives
\begin{equation}
A\star B
=AB+\frac{i}{2}\{A,B\}_{\rm PB}+\mathcal O(\partial_X^2),
\end{equation}
where
\begin{equation}
\{A,B\}_{\rm PB}
=\partial_XA\cdot\partial_pB
-\partial_pA\cdot\partial_XB.
\end{equation}
This is the formal origin of semiclassical drift terms and the controlled gradient expansion used in relativistic kinetic theory.

\section{Clifford decomposition and the axial kinetic sector}

A Dirac Wigner function can be decomposed as
\begin{equation}
\boxed{
W
=\frac14\left[
\mathcal F
+i\gamma^5\mathcal P
+\gamma^\mu\mathcal V_\mu
+\gamma^\mu\gamma^5\mathcal A_\mu
+\frac12\sigma^{\mu\nu}\mathcal S_{\mu\nu}
\right].
}
\label{eq:clifford}
\end{equation}
The coefficient functions are obtained by gamma traces. For example,
\begin{equation}
\mathcal A^\mu
=\Trc(\gamma^\mu\gamma^5W).
\end{equation}

For massive fermions, spin is an independent dynamical degree of freedom rather than being completely locked to momentum. The axial kinetic theory of Hattori, Hidaka, and Yang provides a relativistic one-particle framework for this sector. Collisional extensions identify spin diffusion and polarization effects in the collision term.

For our project, the main lesson is not to copy a published one-body equation verbatim but to preserve its structural ingredients when constructing the pair problem:

\begin{itemize}[leftmargin=1.5em]
\item Lorentz-covariant axial information,
\item gauge-covariant derivatives and color structure,
\item collision/self-energy terms,
\item and the correct momentum-dependent spin projection.
\end{itemize}

\section{Why the pair problem requires a connected four-point function}

A one-particle Wigner function cannot determine a genuine two-particle spin correlation. The relevant object is schematically
\begin{equation}
G^{<(2)}(1,2;1',2')
=\langle
\bar\psi(1')\psi(1)
\bar\psi(2')\psi(2)
\rangle.
\label{eq:fourpoint}
\end{equation}
Its connected piece is
\begin{equation}
G_c^{(2)}
=G^{(2)}-G^{(1)}G^{(1)}
\end{equation}
with the appropriate fermionic exchange terms included according to the species and ordering.

For an $s\bar s$ pair the two legs are distinguishable by particle/antiparticle quantum numbers, but color and spinor indices remain nontrivial.

A double Wigner transform introduces
\begin{equation}
(X_1,p_1;X_2,p_2).
\end{equation}
The axial--axial pair component is obtained by tracing with
\begin{equation}
\gamma^\mu\gamma^5\otimes\gamma^\nu\gamma^5.
\end{equation}
Finally Eq.~\eqref{eq:tetradprojection} converts the Lorentz tensor into the physical $C_{ab}$.

This chain can be summarized as
\begin{equation}
\boxed{
G_c^{(2)}
\longrightarrow
\Ccal^{\mu\nu}_{AA}
\longrightarrow
C_{ab}
\longrightarrow
P.
}
\label{eq:chain4point}
\end{equation}

The outstanding derivational challenge is therefore not merely a more complicated one-particle Wigner equation; it is a kinetic equation for the connected four-point function or an equivalent Bethe-Salpeter/BBGKY object.

\section{Expected structure of the pair kinetic equation}

Without claiming a completed derivation, the exact hierarchy suggests a schematic structure
\begin{equation}
\boxed{
\Dcal_{12}\Ccal_{AA}
=
\Kcal_{\rm drift}[\Ccal_{AA}]
+\Kcal_{\rm coh}[\Ccal_{AA}]
+\Kcal_{\rm coll}[\Ccal_{AA}]
+\Scal_{\rm pair}.
}
\label{eq:schematicpair}
\end{equation}
Here

\begin{itemize}[leftmargin=1.5em]
\item $\Kcal_{\rm drift}$ propagates the two phase-space legs;
\item $\Kcal_{\rm coh}$ contains coherent spin precession and spin-orbit structures;
\item $\Kcal_{\rm coll}$ contains irreversible scattering/dephasing generated by self-energies;
\item $\Scal_{\rm pair}$ is a genuine correlated source associated with pair production or connected gain terms.
\end{itemize}

The white-noise relative-field equation should emerge only after reducing Eq.~\eqref{eq:schematicpair} under strong assumptions: Gaussian bath, short memory, weak backreaction, and projection onto a low-dimensional spin subspace.

This provides a falsifiable consistency condition. A microscopic derivation that does not reduce to common-mode cancellation in the appropriate symmetric limit would be inconsistent with the exact spin algebra derived earlier.

\section{Source--propagator separation and double counting}

Let the pair source be time dependent. A linear tensor equation may be written
\begin{equation}
\dot{\bm C}(t)=M(t)\bm C(t)+\bm G(t),
\end{equation}
where $\bm C$ denotes a vector containing independent tensor components. The exact Duhamel solution is
\begin{equation}
\boxed{
\bm C(t_f)
=U(t_f,t_0)\bm C(t_0)
+\int_{t_0}^{t_f}dt'\,U(t_f,t')\bm G(t'),
}
\label{eq:duhamel}
\end{equation}
with
\begin{equation}
U(t,t')=\mathcal T\exp\left[\int_{t'}^t du\,M(u)\right].
\end{equation}

This equation clarifies a recurring conceptual issue. If early color fields generate the pair correlation during production, that contribution belongs in $\bm G$ or in the source initial condition. Later unital decoherence belongs in $U$. Treating the same color-field correlator as both an additive source and a multiplicative decoherence term without a microscopic derivation risks double counting.

If the source switches off at a time $t_s$, Eq.~\eqref{eq:duhamel} reduces to
\begin{equation}
\bm C(t_f)=U(t_f,t_s)\bm C_{\rm src},
\end{equation}
which is the regime where a source-times-survival factorization is justified.

\section{Hadronization and spin-transfer map}

\subsection{General tensor map}

At hadronization the quark tensor is not directly observed. For a given production/feed-down channel $a$ on the $\Lambda$ side and $b$ on the $\bar\Lambda$ side, write a linear spin-transfer map
\begin{equation}
C^{\Lambda\bar\Lambda,(ab)}_{ij}
=(T_a)_{ik}
C^{s\bar s,(ab)}_{k\ell}
(\bar T_b)_{j\ell}
+C^{\rm conn,had,(ab)}_{ij}.
\label{eq:hadmap}
\end{equation}
The second term allows hadronization itself to create connected spin structure not reducible to independent one-body transfer.

After summing source fractions,
\begin{equation}
\boxed{
C^{\Lambda\bar\Lambda}_{ij}
=\sum_{ab}R_{ab}
(T_a)_{ik}
C^{s\bar s,(ab)}_{k\ell}
(\bar T_b)_{j\ell}
+\sum_{ab}R_{ab}C^{\rm conn,had,(ab)}_{ij}.
}
\label{eq:hadfull}
\end{equation}

A single scalar ``feed-down dilution'' is only valid if all $T_a$ are proportional to the same rotation-independent scalar and if $R_{ab}$ is independent of the kinematic variable being studied. Neither condition is guaranteed.

\subsection{Phenomenological compressed form}

For orientation, one may compress the observable to
\begin{equation}
\boxed{
P_{\Lambda\bar\Lambda}^{\rm obs}
=D_{\rm had}\,
f_{\rm same}\,
P_{s\bar s}^{\rm src}\,
S_{\rm rel}
+P_{\rm conn}^{\rm had/resc},
}
\label{eq:phenomcompact}
\end{equation}
where

\begin{itemize}[leftmargin=1.5em]
\item $D_{\rm had}$ summarizes spin transfer and decay effects,
\item $f_{\rm same}$ is the fraction of reconstructed pairs that retain a common correlated ancestor/source,
\item $S_{\rm rel}$ is the survival factor generated by relative-noise dynamics,
\item $P_{\rm conn}^{\rm had/resc}$ contains correlation created later by hadronization or rescattering.
\end{itemize}

Equation~\eqref{eq:phenomcompact} is useful for bookkeeping but should not be mistaken for a first-principles result. It is the scalar shadow of Eq.~\eqref{eq:hadfull}.

\section{STAR versus CMS: phenomenological constraints}

The STAR $pp$ result at $\sqrt{s}=200$ GeV shows a clearly positive short-range $\Lambda\bar\Lambda$ correlation and a substantial reduction at large separation. CMS at $\sqrt{s}=13$ TeV finds the $\Lambda\bar\Lambda$ correlation consistent with zero within uncertainties, while same-sign pairs are negative at small $\Delta R$.

The difference suggests several nonexclusive possibilities:

\begin{enumerate}[leftmargin=1.6em]
\item The fraction of pairs inheriting a common nonperturbative $s\bar s$ source may decrease with collision energy because perturbative production and independent fragmentation become more important.
\item The hadronization environment may have a different spatial/temporal correlation structure at LHC energies, changing $S_{\rm rel}$.
\item Source composition and feed-down fractions may differ strongly between the STAR and CMS acceptances.
\item The mapping between coordinate-space pair separation during confinement and the measured $\Delta R$ may be energy dependent.
\item Additional same-sign mechanisms, including fragmentation and quantum-statistical/correlation effects, may contribute differently to $\Lambda\Lambda$ and $\bar\Lambda\bar\Lambda$.
\end{enumerate}

The current data therefore do not support a universal one-parameter decoherence curve that can simply be transported from RHIC to the LHC. Energy dependence must be built into the source and transport sectors separately.

\section{Numerical estimates and limiting models}

\subsection{Integrated relative-noise exponent}

For a time-dependent separation $r(t)$,
\begin{equation}
P(t_h)
=P(t_0)
\exp\left[-\int_{t_0}^{t_h}dt\,\Gamma_{\rm rel}(r(t))\right].
\end{equation}
With Eq.~\eqref{eq:Gammarel},
\begin{equation}
D_{\rm rel}
=4\int dt\,[D(0)-D(r(t))].
\end{equation}
The observed correlation measures this integrated quantity, not $D(0)$ or a formation time separately. Therefore a fitted ``decoherence time'' is generally a convolution of microscopic noise strength, spatial cross correlation, and trajectory history.

\subsection{Ballistic toy trajectory}

As a diagnostic, take
\begin{equation}
r(t)=v_{\rm rel}t\,\Delta R
\end{equation}
and a Gaussian spatial correlator
\begin{equation}
D(r)=D_0e^{-r^2/(2\ell_c^2)}.
\end{equation}
Then
\begin{align}
D_{\rm rel}(\Delta R)
&=4D_0\int_0^{\tau_h}dt\,
\left[1-e^{-v_{\rm rel}^2t^2\Delta R^2/(2\ell_c^2)}\right].
\end{align}
Let
\begin{equation}
\nu=\frac{v_{\rm rel}\tau_h\Delta R}{\ell_c}.
\end{equation}
Using
\begin{equation}
\int_0^{\tau_h}dt\,e^{-a^2t^2/2}
=\tau_h\sqrt{\frac\pi2}\frac{1}{\nu}
\operatorname{erf}\left(\frac{\nu}{\sqrt2}\right),
\end{equation}
we obtain
\begin{equation}
\boxed{
D_{\rm rel}(\Delta R)
=4D_0\tau_h
\left[
1-\sqrt{\frac\pi2}\frac1\nu
\operatorname{erf}\left(\frac{\nu}{\sqrt2}\right)
\right].
}
\label{eq:gaussDR}
\end{equation}
At small $\nu$,
\begin{equation}
\operatorname{erf}(\nu/\sqrt2)
=\sqrt{\frac2\pi}\left(\nu-\frac{\nu^3}{6}+\cdots\right),
\end{equation}
so
\begin{equation}
D_{\rm rel}(\Delta R)
\simeq\frac{2}{3}D_0\tau_h\nu^2+\cdots.
\end{equation}
Thus
\begin{equation}
1-\frac{P(\Delta R)}{P(0)}\propto(\Delta R)^2
\end{equation}
near zero separation in this smooth model.

This result is not a prediction until the trajectory ansatz is justified, but it illustrates how a microscopic correlation length can appear in the measured angular dependence.

\section{Positivity and density-matrix consistency checks}

Any phenomenological tensor must correspond to a positive density matrix. For a Bell-diagonal state
\begin{equation}
\rho
=\frac14\left[
\Id\otimes\Id+
C_x\sigma_x\otimes\sigma_x+
C_y\sigma_y\otimes\sigma_y+
C_z\sigma_z\otimes\sigma_z
\right],
\end{equation}
the four eigenvalues are
\begin{align}
\lambda_1&=\frac14(1-C_x-C_y-C_z),\\
\lambda_2&=\frac14(1-C_x+C_y+C_z),\\
\lambda_3&=\frac14(1+C_x-C_y+C_z),\\
\lambda_4&=\frac14(1+C_x+C_y-C_z),
\end{align}
up to the chosen Bell-state ordering. Positivity requires every $\lambda_i\ge0$.

For an isotropic tensor $C_x=C_y=C_z=P$,
\begin{equation}
\lambda_S=\frac14(1-3P),
\qquad
\lambda_T=\frac14(1+P)
\end{equation}
for the singlet and each triplet component, respectively. Thus
\begin{equation}
-1\le P\le\frac13.
\end{equation}
The upper bound $P=1/3$ corresponds to a state with zero singlet weight. Any fit producing $P>1/3$ would be unphysical within the two-spin normalization used here.

These positivity checks should be enforced directly in future likelihood fits rather than applied only after fitting.

\section{Tensor observables beyond the scalar}

If an approximate symmetry axis exists, write
\begin{equation}
C_{ij}=C_T(\delta_{ij}-n_in_j)+C_Ln_in_j.
\end{equation}
Then
\begin{equation}
P=\frac{2C_T+C_L}{3}.
\end{equation}
One scalar cannot determine two functions $C_T$ and $C_L$. This is a fundamental information loss, not merely a statistical limitation.

A useful anisotropy observable is
\begin{equation}
A_C=C_T-C_L.
\end{equation}
Another is
\begin{equation}
R_{LT}=\ln\left|\frac{C_L}{C_T}\right|.
\end{equation}
In a source-dominated multiplicative regime,
\begin{equation}
C_T=C_T^{\rm src}e^{-D_T},
\qquad
C_L=C_L^{\rm src}e^{-D_L},
\end{equation}
so
\begin{equation}
R_{LT}
=\ln\left|\frac{C_L^{\rm src}}{C_T^{\rm src}}\right|
-(D_L-D_T).
\end{equation}
If the source ratio is independently constrained, $R_{LT}$ directly measures anisotropic decoherence. This would discriminate far more strongly among microscopic models than the trace alone.

\section{Corrections and clarifications relative to earlier stages}

Several conceptual points were sharpened on 23 August.

\subsection{Connected versus full correlation}

Earlier scalar formulas often implicitly assumed $\bm p=\bar{\bm p}=0$. The correct general object is Eq.~\eqref{eq:Cconnected}. This matters for heavy-ion collisions where one-body hyperon polarization may be nonzero.

\subsection{Relativistic projection}

A naive spatial trace of an axial--axial Lorentz tensor is not generally the physical spin scalar. Momentum-dependent tetrads are required, Eq.~\eqref{eq:relP}.

\subsection{Relative-noise rate}

The important rate is not a generic local noise strength but the difference between local and cross noise. In the symmetric white-noise limit the exact combination is Eq.~\eqref{eq:Gammarel}.

\subsection{Generation versus decoherence}

A Hermitian Gaussian stochastic field generates a unital channel. It cannot create spin correlation from the maximally mixed state. Correlation generated by early color fields belongs to the source/four-point sector, not to the same late-time unital damping term.

\subsection{Perturbative spin-flip estimates}

Long perturbative strange-quark spin-equilibration times do not invalidate nonperturbative hadronization decoherence; they only show that ordinary weak-coupling spin-flip scattering is unlikely to be fast enough by itself.

\section{Concrete microscopic program emerging from the day}

The research program can now be stated more sharply.

\subsection{Step 1: source calculation}

Construct the correlated $s\bar s$ source density matrix or four-point function for the relevant production mechanisms:

\begin{itemize}[leftmargin=1.5em]
\item vacuum/string-breaking-like production,
\item perturbative $g\to s\bar s$ production,
\item and, in heavy-ion collisions, glasma/Schwinger production.
\end{itemize}

These mechanisms need not have the same initial tensor.

\subsection{Step 2: connected two-particle kinetic equation}

Derive a gauge-covariant equation for $G_c^{(2)}$ or its axial--axial projection. The collision kernel should retain explicit spinor and color indices long enough to identify which terms are common mode and which are relative mode.

\subsection{Step 3: controlled reduction}

Under stated approximations, reduce the microscopic equation to an effective form and verify that it reproduces
\begin{equation}
\Gamma_{\rm rel}\propto D_{11}+D_{22}-D_{12}-D_{21}
\end{equation}
in the classical Gaussian limit.

\subsection{Step 4: trajectory and geometry model}

Relate the space-time separation entering the QCD correlator to measured pair kinematics. This may require a transport/event-generator interface rather than an analytic $r\propto t\Delta R$ ansatz.

\subsection{Step 5: hadronization and feed-down}

Implement Eq.~\eqref{eq:hadfull} channel by channel for primary $\Lambda$, $\Sigma^0$, $\Xi$, and other relevant parents, keeping tensor rather than scalar transfer whenever possible.

\subsection{Step 6: simultaneous RHIC/LHC test}

A credible model should explain why the positive short-range $\Lambda\bar\Lambda$ signal is prominent at STAR but not statistically significant at CMS, rather than fitting only one collision energy.

\section{Highest-priority derivation for the next research day}

The single highest-value technical step is to write the connected four-point Kadanoff--Baym equation with explicit indices. Let
\begin{equation}
G_{\alpha\beta;\gamma\delta}^{(2)<}(1,2;1',2')
\end{equation}
carry spinor indices, with color indices suppressed here for readability. One should act with the Dirac operator on leg 1 and separately on leg 2, yielding two coupled equations. After subtracting products of one-body equations, the connected piece will have schematic terms
\begin{equation}
\left(G^{-1}_{0,1}-\Sigma_1\right)\circ G_c^{(2)}
=\Scal_{12}+\Kcal_{12}[G_c^{(2)}],
\end{equation}
and an analogous equation for leg 2.

The immediate tasks are:

\begin{enumerate}[leftmargin=1.6em]
\item keep color indices and introduce gauge links before Wigner transformation;
\item perform the double Wigner transform;
\item project the connected function with $\gamma^\mu\gamma^5\otimes\gamma^\nu\gamma^5$;
\item contract with the spin tetrads;
\item identify the tensor structures that reduce to common and relative noise;
\item project further onto singlet/triplet sectors if the scalar is the first target;
\item derive the weak-coupling/short-memory limit and compare it with Eq.~\eqref{eq:Gammarel}.
\end{enumerate}

This route would turn the current phenomenological insight into a genuinely microscopic transport theory.

\section{Summary of validated conclusions from 23 August}

The full-day analysis supports the following conclusions.

\begin{enumerate}[leftmargin=1.6em]
\item The correct object is a connected two-particle spin tensor, not merely two independent one-body polarizations.
\item The exact BBGKY reduction shows that the pair equation couples to three-body correlations; every closed kinetic equation therefore contains a nontrivial approximation.
\item Relativistic spin correlations must be projected from the axial tensor with momentum-dependent spin tetrads.
\item Independent one-body depolarization damps the pair mode with the sum of the two leg rates.
\item Perfectly common stochastic spin fields leave $\bm\sigma_1\cdot\bm\sigma_2$ invariant exactly.
\item In the isotropic Gaussian white-noise limit, the scalar damping rate is controlled by the relative covariance, $4[D(0)-D(r)]$ for equivalent trajectories.
\item Finite bath memory produces a quadratic short-time decay and an exponential long-time limit while preserving the same relative-noise combination.
\item A gauge-covariant microscopic interpretation requires Wilson-line-connected QCD field-strength correlators.
\item For strange quarks, a heavy-quark Pauli coefficient is not a controlled quantitative matching prescription; a relativistic axial-kinetic derivation is preferable.
\item Hermitian Gaussian noise is unital and cannot create pair correlation from the maximally mixed state; source generation and later decoherence must be separated.
\item The pair problem requires a connected four-point Kadanoff--Baym/Wigner object and an axial--axial projection.
\item Hadronization, feed-down, source fractions, and late connected effects are indispensable for comparison with $\Lambda\bar\Lambda$ data.
\item The STAR/CMS contrast indicates nontrivial energy dependence and argues against a universal one-parameter decoherence model.
\end{enumerate}

\section*{References}
\begin{enumerate}[leftmargin=1.6em]
\item STAR Collaboration, ``Measuring spin correlation between quarks during QCD confinement,'' \emph{Nature} \textbf{650}, 65--71 (2026), arXiv:2506.05499.
\item CMS Collaboration, ``Measurements of $\Lambda$ hyperon spin correlations in proton-proton and proton-lead collisions,'' CMS-PAS-HIN-26-002 (17 May 2026).
\item J. I. Kapusta, E. Rrapaj and S. Rudaz, ``Relaxation time for strange quark spin in rotating quark-gluon plasma,'' \emph{Phys. Rev. C} \textbf{101}, 024907 (2020).
\item J. I. Kapusta, E. Rrapaj and S. Rudaz, ``Spin versus helicity equilibration times and Lagrangian for strange quarks in rotating quark-gluon plasma,'' \emph{Phys. Rev. C} \textbf{102}, 064911 (2020).
\item K. Hattori, Y. Hidaka and D.-L. Yang, ``Axial Kinetic Theory and Spin Transport for Fermions with Arbitrary Mass,'' \emph{Phys. Rev. D} \textbf{100}, 096011 (2019), arXiv:1903.01653.
\item D.-L. Yang, K. Hattori and Y. Hidaka, ``Effective quantum kinetic theory for spin transport of fermions with collisional effects,'' arXiv:2002.02612.
\item D.-L. Yang, ``Quantum kinetic theory for spin transport of quarks with background chromo-electromagnetic fields,'' \emph{JHEP} \textbf{06} (2022) 140, arXiv:2112.14392.
\item A. Kumar, B. M\"uller and D.-L. Yang, ``Spin polarization and correlation of quarks from the glasma,'' \emph{Phys. Rev. D} \textbf{107}, 076025 (2023), arXiv:2212.13354.
\item J. Barata, I. Cunt\'in, E. Rico and B. Wu, ``$\Lambda\bar\Lambda$ spin correlations in high-energy collisions from quantum channels: an open quantum system view of hadronization,'' arXiv:2606.30737 (2026).
\end{enumerate}

\end{document}
